% Generated by report_gen.py (ai-ee v1) - do not hand-edit.
\documentclass[11pt,a4paper]{article}
\usepackage[margin=2.2cm]{geometry}
\usepackage{graphicx}
\usepackage{booktabs}
\usepackage{longtable}
\usepackage{pdfpages}
\makeatletter
\@ifundefined{KV@Gin@artifact}{\define@key{Gin}{artifact}[]{}}{}
\makeatother
\usepackage{xcolor}
\usepackage[hidelinks]{hyperref}
\setcounter{tocdepth}{1}
\begin{document}
\sloppy
\begin{center}
{\LARGE\bfseries lumina-strobe --- Design Document}\\[6pt]
{\large ai-ee v1 pipeline}\\[2pt]
generated 2026-07-28 19:40:46
\end{center}
\tableofcontents

\section{Board and Run Metadata}
{\small
\begin{longtable}{lp{11cm}}
\toprule
\textbf{Field} & \textbf{Value} \\
\midrule
\endhead
board & lumina-strobe \\
workspace & \texttt{boards/lumina-strobe} \\
phase & P2 \\
gate status & no gates recorded yet \\
state created & 2026-07-28T14:59:51 \\
state updated & 2026-07-28T19:39:41 \\
\bottomrule
\end{longtable}
}
\section{Overview}
\section{Requirements}
\subsection*{Requirements: LUM-DTR-STROBE-A (LUMINA strobe daughter board)}
Sources, in precedence order (later overrides earlier where they conflict):

1. \texttt{brief/00-lumina-system-context.md} - shared system context. 2. \texttt{brief/01-carrier-board-brief.md} - the mating carrier. CONTEXT ONLY, not this board. 3. \texttt{brief/02-strobe-daughter-brief.md} - this board's brief. 4. \texttt{brief/05-lumina-closed-decisions.md} - BINDING. Supersedes \texttt{00} section 6 and corrects several figures in \texttt{00} section 5 and \texttt{02} section 3. 5. \texttt{brief/06-connector-icd.md} - the frozen expansion connector ICD. A HARD input. Never redefined by this board. Per its own preamble, where a number appears in more than one place the ICD is the one that governs.

Unstated low-risk items are marked \texttt{ASSUMED:} inline. Section 9 holds every design-changing or safety-relevant unknown. D-01, D-02 and D-03 are CLOSED and are recorded here as settled requirements, not questions.

---

\subsubsection*{1. Function}
A fixture-specific daughter board that stacks above the universal LUMINA carrier (LUM-CAR-A) on 11.0 mm standoffs and mates through the frozen 38-position expansion connector. It contains the entire light engine for a strobe fixture: a local energy store on the 48 V raw PD rail, an inrush limiter, a mandatory bleed path, a drive stage that dumps that store into a series white LED string in short high-current bursts with hard optical edges, sense outputs (bank voltage, LED temperature) back to the carrier's ADCs, an over-temperature shutdown that acts independently of firmware, and the board-ID mechanism that lets one firmware image recognise this as a strobe.

It carries no MCU, no network interface, no power conversion from mains, and no external connector of any kind. Strobe pattern generation, the average-energy governor and colour algorithms are host/firmware, not hardware - the hardware's job is to expose what the governor needs and to fail safe without it.

Quantity: 4-6 boards, part of a first build of 8-12 fixtures.

\textbf{The board exists to solve one problem:} turn a \textasciitilde{}8.5 W sustained power budget into short, violent bursts of light with sub-millisecond edges, and degrade gracefully (dimmer flashes, never missed flashes, never a reset) when the music asks for more than the rail can deliver.

---

\subsubsection*{2. Interfaces}
All external electrical connection is through two board-to-board sockets on the \textbf{bottom} side of this board. \textbf{No other connector may be added} (ICD s9): no barrel jack, no DMX, no second Ethernet, no USB. The only permitted additional wiring is internal to the enclosure (the LED string - see open question 4).

\paragraph*{2.1 J3 - POWER block}
\begin{itemize}
\item Part: {*}\emph{CONNFLY DS1023-2}7SF11{*}{*}, 14 positions, 2.54 mm, 600 V, 3.0 A/contact, 8.5 mm body.
\item Orientation: \textbf{faces downward} - reverse-mounted THT on the bottom side, or a bottom-side
SMD equivalent (ICD s7.3). See open question 6.
\item Position 1 silkscreen-marked with a triangle.
\item Nominal body extent \textbf{(14, 68) - (34, 78)}; position 1 at (15.3, 69.3), long axis along +x.
\textbf{Provisional until the end of the carrier's P6} (ICD s7.2) - see open question 9.
\end{itemize}
\begin{quote}\small\ttfamily\raggedright
\textbar{} Col \textbar{} Row A \textbar{} net \textbar{} Row B \textbar{} net \textbar{}\\
\textbar{}---\textbar{}---\textbar{}---\textbar{}---\textbar{}---\textbar{}\\
\textbar{} 1 \textbar{} 1 \textbar{} `+48V\_SW` \textbar{} 2 \textbar{} `GND` \textbar{}\\
\textbar{} 2 \textbar{} 3 \textbar{} `+48V\_SW` \textbar{} 4 \textbar{} `GND` \textbar{}\\
\textbar{} 3 \textbar{} 5 \textbar{} `+48V\_SW` \textbar{} 6 \textbar{} `GND` \textbar{}\\
\textbar{} 4 \textbar{} 7 \textbar{} `GND` (guard column) \textbar{} 8 \textbar{} `GND` \textbar{}\\
\textbar{} 5 \textbar{} 9 \textbar{} `+12V` \textbar{} 10 \textbar{} `GND` \textbar{}\\
\textbar{} 6 \textbar{} 11 \textbar{} `+12V` \textbar{} 12 \textbar{} `+3V3` \textbar{}\\
\textbar{} 7 \textbar{} 13 \textbar{} `GND` \textbar{} 14 \textbar{} `+3V3` \textbar{}
\end{quote}
Design properties this map is built to have, each checkable by eye (ICD s3.1): the 48 V group is at one end and bounded by GND on every side; \textbf{column 4 is an all-GND guard column}; every supply pin has an adjacent GND; rail order along the connector is {*}{*}48 -\textgreater{} GND -\textgreater{} 12 -\textgreater{} 3.3{*}{*}, so no single-position mis-seat can put a higher rail on a lower rail's pin.

\paragraph*{2.2 J4 - SIGNAL block}
\begin{itemize}
\item Part: {*}\emph{CONNFLY DS1023-2}12SF11{*}{*}, 24 positions, 2.54 mm, 600 V, 3.0 A/contact, 8.5 mm body.
\item Orientation, marking and provisional-coordinate status: as J3.
\item Nominal body extent (56, 68) - (88, 78); position 1 at (57.3, 69.3), long axis along +x.
\end{itemize}
\begin{quote}\small\ttfamily\raggedright
\textbar{} Col \textbar{} Row A \textbar{} net \textbar{} Row B \textbar{} net \textbar{}\\
\textbar{}---\textbar{}---\textbar{}---\textbar{}---\textbar{}---\textbar{}\\
\textbar{} 1 \textbar{} 1 \textbar{} `PWM0` \textbar{} 2 \textbar{} `PWM1` \textbar{}\\
\textbar{} 2 \textbar{} 3 \textbar{} `GND` \textbar{} 4 \textbar{} `GND` \textbar{}\\
\textbar{} 3 \textbar{} 5 \textbar{} `PWM2` \textbar{} 6 \textbar{} `PWM3` \textbar{}\\
\textbar{} 4 \textbar{} 7 \textbar{} `PWM4` \textbar{} 8 \textbar{} `PWM5` \textbar{}\\
\textbar{} 5 \textbar{} 9 \textbar{} `GND` \textbar{} 10 \textbar{} `GND` \textbar{}\\
\textbar{} 6 \textbar{} 11 \textbar{} `PWM6` \textbar{} 12 \textbar{} `PWM7` \textbar{}\\
\textbar{} 7 \textbar{} 13 \textbar{} `GND` \textbar{} 14 \textbar{} `DSPI\_SCK` \textbar{}\\
\textbar{} 8 \textbar{} 15 \textbar{} `DSPI\_MOSI` \textbar{} 16 \textbar{} `DSPI\_MISO` \textbar{}\\
\textbar{} 9 \textbar{} 17 \textbar{} `DSPI\_CSn` \textbar{} 18 \textbar{} `I2C\_SCL` \textbar{}\\
\textbar{} 10 \textbar{} 19 \textbar{} `I2C\_SDA` \textbar{} 20 \textbar{} `ADC0` \textbar{}\\
\textbar{} 11 \textbar{} 21 \textbar{} `ADC1` \textbar{} 22 \textbar{} `ID\_ADC` \textbar{}\\
\textbar{} 12 \textbar{} 23 \textbar{} `ENABLE` \textbar{} 24 \textbar{} `FAULT` \textbar{}
\end{quote}
\textbf{No 48 V exists anywhere on J4.}

\paragraph*{2.3 Signal electrical contract (ICD s3.3) - binding}
\begin{quote}\small\ttfamily\raggedright
\textbar{} Signal \textbar{} Direction \textbar{} Rules this board must honour \textbar{}\\
\textbar{}---\textbar{}---\textbar{}---\textbar{}\\
\textbar{} `PWM0..7` \textbar{} carrier -\textgreater{} daughter \textbar{} 3.3 V push-pull CMOS. Default {*}{*}13-bit at 9.766 kHz{*}{*}. `PWM0-3` = LEDC timer 0, `PWM4-7` = LEDC timer 1. Channels on one timer share frequency AND resolution (CAR-REQ-11). Max 8 channels total (CAR-REQ-10) \textbar{}\\
\textbar{} `DSPI\_{*}` \textbar{} bidirectional \textbar{} SPI mode 0, MSB first, \textless{}= 26 MHz. This board drives no line except MISO, and only while `DSPI\_CSn` is low. Shared bus, one CS \textbar{}\\
\textbar{} `I2C\_SCL/SDA` \textbar{} bidirectional \textbar{} Open drain, 400 kHz. {*}{*}Pull-ups are on the carrier (4.7 k). This board must NOT fit its own{*}{*} \textbar{}\\
\textbar{} `ADC0`, `ADC1` \textbar{} daughter -\textgreater{} carrier \textbar{} Analogue 0-3.3 V. {*}{*}Source impedance \textless{}= 10 kohm{*}{*} \textbar{}\\
\textbar{} `ID\_ADC` \textbar{} daughter -\textgreater{} carrier \textbar{} Carrier fits the top leg (10 k to +3V3); {*}{*}this board fits the bottom leg to GND{*}{*}. Code allocated by the carrier owner - see open question 8 \textbar{}\\
\textbar{} `ENABLE` \textbar{} carrier -\textgreater{} daughter \textbar{} Push-pull CMOS, {*}{*}active HIGH{*}{*}. This board fits its own {*}{*}100 kohm pull-down{*}{*}, gates {*}{*}every{*}{*} output stage with it (including the cap-bank charge path), and {*}{*}never latches it locally{*}{*} \textbar{}\\
\textbar{} `FAULT` \textbar{} daughter -\textgreater{} carrier \textbar{} {*}{*}Open drain, active low{*}{*}, 10 k pull-up on the carrier, wire-OR'd with the carrier's eFuse fault. {*}{*}Never drive it high{*}{*} \textbar{}
\end{quote}
\paragraph*{2.4 Carrier-facing functional requirements (from \texttt{02} section 6)}
\begin{itemize}
\item \textbf{STR-REQ-17} - PWM channel count within the 8-channel ceiling. White-only needs 1-2;
RGBW needs 4-5. Resolved by open question 1.
\item \textbf{STR-REQ-18} - report bank voltage to a carrier ADC pin (feeds the governor, STR-REQ-06).
\item \textbf{STR-REQ-19} - report LED thermistor temperature to a carrier ADC or I2C pin.
\item \textbf{STR-REQ-20} - assert FAULT on over-temperature and shut down {*}{*}independently of
firmware{*}{*}. Do not rely on the network or the MCU to prevent a thermal event.
\item \textbf{STR-REQ-21} - honour the fail-safe ENABLE: outputs off with ENABLE de-asserted,
including during MCU reset and firmware update.
\item \textbf{STR-REQ-22} - populate the board-ID mechanism (CAR-REQ-07).
\end{itemize}
\texttt{ASSUMED:} bank voltage on \texttt{ADC0}, LED thermistor on \texttt{ADC1}, board ID on \texttt{ID\_ADC}. If open question 1 closes as RGBW and per-channel sense is wanted, the sense budget re-opens and the extra channels must go on I2C.

\texttt{ASSUMED:} no SPI device on this board; \texttt{DSPI\_{*}} left unconnected at the socket. Revisit only if open question 1 closes as RGBW and the PWM budget forces an external LED driver.

\paragraph*{2.5 LED string interface}
\begin{itemize}
\item Series LED string driven from the 48 V bank (STR-REQ-12): \textasciitilde{}2.6 A at roughly 38 V string
voltage, versus \textasciitilde{}8.3 A at 12 V. This is a closed consequence of D-02.
\item Whether the emitters sit on this PCB or on a remote heatsink wired to it is \textbf{not settled}
\begin{itemize}
\item see open question 4. Either way the string, its wiring and its heatsink are at PoE
potential (section 8).
\end{itemize}
\end{itemize}
---

\subsubsection*{3. Power}
\paragraph*{3.1 What the connector delivers - ICD s6.2, binding}
\begin{quote}\small\ttfamily\raggedright
\textbar{} Rail \textbar{} Sustained (af) \textbar{} Sustained (at) \textbar{} Peak, ms-scale, from local bulk \textbar{} Hardware fault ceiling \textbar{} Connector pin capacity \textbar{}\\
\textbar{}---\textbar{}---\textbar{}---\textbar{}---\textbar{}---\textbar{}---\textbar{}\\
\textbar{} `+48V\_SW` \textbar{} {*}{*}0.25 A{*}{*} \textbar{} 0.50 A \textbar{} 1.0 A (eFuse limit, until it latches) \textbar{} 1.0 A, {*}{*}latch off{*}{*} \textbar{} 5.40 A \textbar{}\\
\textbar{} `+12V` \textbar{} 0.75 A \textbar{} 1.25 A \textbar{} 2.0 A \textbar{} 2.0 A converter, OCP \textbar{} 3.60 A \textbar{}\\
\textbar{} `+3V3` \textbar{} 0.25 A \textbar{} 0.25 A \textbar{} 0.50 A \textbar{} 1.0 A converter \textbar{} 3.60 A \textbar{}\\
\textbar{} {*}{*}TOTAL, all three rails{*}{*} \textbar{} {*}{*}8.5 W{*}{*} \textbar{} {*}{*}18.5 W{*}{*} \textbar{} - \textbar{} PSE overload timer \textasciitilde{}50-75 ms \textbar{} - \textbar{}
\end{quote}
\textbf{The per-rail ceilings do not add up and are not meant to. The TOTAL is what binds.} \texttt{05} states the closed-decision budget as 8.6-9.3 W (af) / 18.7-20.0 W (at) after 2.4/3.7 W carrier overhead; the ICD's 8.5 W / 18.5 W governs where they differ. Design against \textbf{af}.

\textbf{Superseded, do not use:} the "\textasciitilde{}10 W regulated / \textasciitilde{}8.5 W to the light engine" chain in \texttt{00} section 5.1 and \texttt{02} section 3.1 derives from a Skyworks Si3402-B PD part this system does not use. The carrier uses a TPS2378-class PD controller plus a 100 V buck (\texttt{05}).

\textbf{D-01 closed:} 802.3\textbf{at} power stage, \textbf{af} classification for the first build (resistor change + a PoE+ switch upgrades it, no respin). \textbf{Daughter budgets are designed against af.} Nothing on this board may become the part that pins the fixture to Type 1.

\textbf{ICD s6.3:} take power on \texttt{+48V\_SW}, not \texttt{+12V}, wherever possible - it skips the 48-\textgreater{}12 conversion and is worth 0.67 W (af) / 1.30 W (at) of extra delivered power, and it removes the carrier's only real hot spot.

\texttt{ASSUMED:} this board's own housekeeping (gate drive bias, comparator, thermistor bias) is a small fraction of the budget and is taken from \texttt{+3V3} and/or \texttt{+12V}; all LED energy comes from the 48 V bank.

\paragraph*{3.2 Energy store - D-02 closed}
\textbf{The energy store runs off the 48 V raw rail from the connector.} Not 33,000 uF at 12 V.

\begin{quote}\small\ttfamily\raggedright
\textbar{} Item \textbar{} Value \textbar{} Source \textbar{}\\
\textbar{}---\textbar{}---\textbar{}---\textbar{}\\
\textbar{} Bank capacitance \textbar{} {*}{*}\textasciitilde{}2,800 uF{*}{*} \textbar{} D-02 closed, ICD s6.4 \textbar{}\\
\textbar{} Capacitor voltage rating \textbar{} {*}{*}\textgreater{}= 100 V{*}{*} (63 V is NOT enough at 57 V worst case once ceramic DC-bias derating applies) \textbar{} ICD s5.4 \textbar{}\\
\textbar{} Energy over the 48 -\textgreater{} 40 V window \textbar{} {*}{*}0.99 J{*}{*} \textbar{} ICD s6.4 \textbar{}\\
\textbar{} Energy at full 0 -\textgreater{} 48 V charge \textbar{} {*}{*}3.23 J{*}{*} \textbar{} ICD s6.4 \textbar{}\\
\textbar{} Max full-window flash rate the rail sustains \textbar{} {*}{*}8.6 Hz (af){*}{*} / 18.8 Hz (at) \textbar{} ICD s6.4 \textbar{}\\
\textbar{} Energy per flash at SYS-REQ-03's 25 Hz ceiling \textbar{} {*}{*}0.34 J (af){*}{*} / 0.74 J (at) \textbar{} ICD s6.4 \textbar{}\\
\textbar{} Bank droop per flash at 25 Hz \textbar{} 48 -\textgreater{} 45.4 V (af) / 48 -\textgreater{} 42.1 V (at) \textbar{} ICD s6.4 \textbar{}\\
\textbar{} Cold-start charge time at the ICD sustained limit \textbar{} {*}{*}0.54 s (af){*}{*} / 0.27 s (at) \textbar{} ICD s6.4 \textbar{}
\end{quote}
SYS-REQ-03's 1-25 Hz range is reachable on af, at reduced per-flash energy above \textasciitilde{}8.6 Hz.

\begin{itemize}
\item \textbf{STR-REQ-08} - select capacitors on \textbf{pulse ripple current and ESR}, not capacitance
alone. A high-ESR bank sags harder than the ideal energy calculation predicts and softens
the flash edge. \textbf{Ripple current rating must appear explicitly in the BOM.}
\item \textbf{STR-REQ-09 / CAR-REQ-14} - inrush limiting is \textbf{this board's} responsibility (NTC or
soft-start MOSFET). {*}{*}Size it against the PD's 1.0 A operating current limit, not against
the connector's 5.4 A rating{*}{*} (ICD s8.2). Sizing against the connector is the classic way
to trip the PD's 800 us foldback deglitch and brown out the entire fixture. The carrier's
eFuse dV/dT is deliberately set fast and its limit sits above this board's inrush level so
the two soft-starts do not fight.
\item \textbf{STR-REQ-10 / CAR-REQ-17} - \textbf{bleed path across the bank is mandatory.} The carrier
bleeds its own side through 100 kohm but {*}{*}deliberately fits no series diode on
\texttt{+48V\_SW}{*}{*}, so this board's bleed path is not stranded above the carrier's.
\end{itemize}
\paragraph*{3.3 802.3 compliance clause - ICD s8.3, binding}
IEEE 802.3 caps PD port capacitance at roughly 180 uF; this board's \textasciitilde{}2,800 uF is 15x that. The carrier's 48 V load switch is therefore a \textbf{compliance part}, off through detection, classification, inrush and the 80 ms window, closing only after firmware asserts ENABLE.

1. \textbf{`+48V\_SW` is DEAD at power-up and stays dead for hundreds of milliseconds.} Design for it. 2. {*}{*}This board must provide NO path that energises its bank from \texttt{+12V} or \texttt{+3V3} while \texttt{+48V\_SW} is off.{*}{*} Doing so re-creates the compliance problem behind the switch's back. 3. \textbf{No mate sequencing exists.} A dual-row header has no first-mate/last-mate control, so this board must tolerate 48 V arriving \textbf{before or after} 3.3 V, in either order.

\paragraph*{3.4 The governor, and a derived tension worth reading}
\begin{itemize}
\item \textbf{STR-REQ-06} - firmware implements a leaky-bucket average-energy governor; the hardware's
obligation is to expose bank rail voltage to the carrier ADC.
\item \textbf{STR-REQ-07} - rail sag must degrade \textbf{gracefully}: dimmer flashes, never missed flashes
or a reset. A dropped beat is more visible than a slightly dimmer one.
\end{itemize}
Derived arithmetic from the closed bank size (this is the substance of open question 2):

\begin{quote}\small\ttfamily\raggedright
\textbar{} Flash duration \textbar{} Energy available (0.99 J bank, full window) \textbar{} Implied drive power \textbar{} Implied string current at \textasciitilde{}38 V \textbar{}\\
\textbar{}---\textbar{}---\textbar{}---\textbar{}---\textbar{}\\
\textbar{} 10 ms \textbar{} 0.99 J \textbar{} \textasciitilde{}99 W \textbar{} \textasciitilde{}2.6 A (the closed figure) \textbar{}\\
\textbar{} 50 ms (STR-REQ-03) \textbar{} 0.99 J \textbar{} \textasciitilde{}20 W \textbar{} \textasciitilde{}0.52 A \textbar{}\\
\textbar{} 100 ms (STR-REQ-01) \textbar{} 0.99 J + rail \textbar{} \textasciitilde{}11-15 W \textbar{} \textasciitilde{}0.3-0.4 A \textbar{}\\
\textbar{} 200 ms (STR-REQ-01) \textbar{} 0.99 J + rail \textbar{} \textasciitilde{}7-11 W \textbar{} \textasciitilde{}0.2-0.3 A \textbar{}
\end{quote}
At the closed 2.6 A drive current the bank holds full output for \textbf{\textasciitilde{}8.6 ms} (\texttt{dt = C{*}dV/I = 2800 uF {*} 8 V / 2.6 A}), not the 100-200 ms STR-REQ-01 asks for. The "100 W strobe" in Project Plan v1.0 and the worked example in \texttt{02} section 3.2 are a \textbf{10 ms} operating point. {*}{*}STR-REQ-01's 100-200 ms full-output flash and the closed bank size are not consistent at 100 W{*}{*} - a 150 ms flash lands near 15 W, roughly 6x below the 2.6 A peak the drive stage is built for. This is not a blocker (the drive stage is still sized for 2.6 A peak for short accents) but it fixes what "full output" means optically, and therefore fixes the LED choice. Put to the human as open question 2.

---

\subsubsection*{4. Environment}
\begin{quote}\small\ttfamily\raggedright
\textbar{} Item \textbar{} Value \textbar{} Source \textbar{}\\
\textbar{}---\textbar{}---\textbar{}---\textbar{}\\
\textbar{} Enclosure \textbar{} {*}{*}Plastic / 3D-printed, non-conductive.{*}{*} Sealed (no forced airflow) \textbar{} H1-Q5, ICD s9 \textbar{}\\
\textbar{} Internal air temperature \textbar{} {*}{*}56 C (af) / 69 C (at){*}{*} \textbar{} ICD s7.6 \textbar{}\\
\textbar{} Connector contact rating \textbar{} -40 to +105 C \textbar{} ICD s1 \textbar{}\\
\textbar{} Room \textbar{} \textasciitilde{}5 x 7 x 2.5 m basement/garage, 8-12 fixtures \textbar{} `00` section 1 \textbar{}\\
\textbar{} Ingress / vibration \textbar{} Not stated \textbar{} - \textbar{}
\end{quote}
\begin{itemize}
\item The 56-69 C internal ambient is a \textbf{capacitor selection constraint}: aluminium
electrolytic life halves per 10 C, and the bank is the board's largest and most
lifetime-sensitive part.
\item \textbf{STR-REQ-15} - thermal path sized for the \textbf{sustained} average, with junction
temperature during peak pulses checked against the pulsed derating curve. {*}{*}Both cases must
pass.{*}{*}
\item \textbf{STR-REQ-05} - survive 8-16 bars of maximum-rate flashing, then return to normal.
\texttt{ASSUMED:} at typical tempos this is \textasciitilde{}20-35 s. \textbf{Duration is not actually a free variable}
for the thermal steady state: the rail caps sustained input at 8.5 W (af) regardless of how
long the drop runs, so the steady-state thermal case is bounded by the rail, not by the
bar count.
\item \texttt{ASSUMED:} indoor, dry, benign. No ingress rating, no vibration or shock requirement, no
condensation. Low risk - the fixture shares an enclosure with a PoE carrier already
designed for the same conditions.
\end{itemize}
---

\subsubsection*{5. Size and mounting}
Common LUMINA footprint (ICD s7.1) - \textbf{inherited, not chosen by this board}:

\begin{quote}\small\ttfamily\raggedright
\textbar{} Item \textbar{} Value \textbar{}\\
\textbar{}---\textbar{}---\textbar{}\\
\textbar{} Board outline \textbar{} {*}{*}100.0 x 80.0 mm{*}{*} \textbar{}\\
\textbar{} Corner radius \textbar{} {*}{*}3.0 mm{*}{*}, all four corners \textbar{}\\
\textbar{} Board thickness \textbar{} {*}{*}1.6 mm{*}{*} \textbar{}\\
\textbar{} Mounting holes \textbar{} {*}{*}4x M3 (3.2 mm) at 5 mm inset{*}{*} (a 90 x 70 mm hole rectangle) {*}{*}plus a 5th M3 at (46, 74){*}{*} \textbar{}\\
\textbar{} Coordinate origin \textbar{} board top-left corner, x right, y down \textbar{}\\
\textbar{} Stack arrangement \textbar{} {*}{*}Stacked mezzanine, daughter above the carrier{*}{*} \textbar{}\\
\textbar{} Mated board-to-board height \textbar{} {*}{*}11.0 mm{*}{*} hard-seated, against a positive mechanical stop \textbar{}\\
\textbar{} Standoffs \textbar{} 5x M3 female-female, 11.0 mm \textbar{}\\
\textbar{} Socket orientation \textbar{} {*}{*}Faces downward{*}{*} \textbar{}
\end{quote}
Per MECH-01, \texttt{board\_init.py} needs \texttt{--corner-radius 3} explicitly (default is 0) and \texttt{--mounting-holes 4}; the radius is clamped to the mounting-hole inset (\texttt{margin / 2}), so 3.0 mm works at the default \texttt{--margin 6}. Read \texttt{corner\_radius} and \texttt{worker\_notes} in the board\_init report - do not assume the requested value was honoured. H5 at (46, 74) is added at P4 as a \texttt{MountingHole\_3.2mm\_M3} symbol so it carries a refdes.

Per MECH-02, \textbf{ai-ee has no outline-shrink step - the P5 outline is final.} The outline, radius and hole pattern above are already closed by the ICD, so there is nothing to negotiate here; but the notch below must be in the P5 outline, not retrofitted.

\paragraph*{5.1 The RJ45 notch - HARD requirement from P2 onward}
\textbf{A 30 x 26 mm relief in the TOP edge, region (6, 0) - (36, 26).} The outline rectangle, corner radius and 5-hole pattern are unchanged; only this local relief differs.

Two independent reasons, both load-bearing:

1. The carrier's board-edge magjack is \textasciitilde{}15 mm tall and the stack is 11.0 mm, so the jack protrudes \textasciitilde{}4 mm above this board's underside. Without the notch the boards {*}{*}cannot be forced flat{*}{*}. 2. It is the \textbf{primary reverse-insertion interlock} (CAR-REQ-16). MECH-01's 4x M3 pattern is rotationally symmetric, so a daughter can be bolted down rotated 180 degrees; rotated, the notch lands at the bottom edge and the board presents solid material over the jack. This is a mechanical stop, not a warning.

\paragraph*{5.2 Exclusion zones - ICD s7.6, all board-relative}
\begin{quote}\small\ttfamily\raggedright
\textbar{} Zone \textbar{} Region \textbar{} Requirement \textbar{}\\
\textbar{}---\textbar{}---\textbar{}---\textbar{}\\
\textbar{} {*}{*}DC-DC hot zone{*}{*} \textbar{} {*}{*}(2, 46) - (36, 68){*}{*} \textbar{} {*}{*}No LED drivers and no aluminium electrolytics.{*}{*} The carrier's 48-\textgreater{}12 converter dissipates up to 1.25 W directly below in a sealed box. This is a vertical keepout, not an in-plane separation rule \textbar{}\\
\textbar{} {*}{*}Antenna column{*}{*} \textbar{} {*}{*}(88, 25) - (100, 55){*}{*} \textbar{} {*}{*}No copper on ANY layer{*}{*} (no traces, no pour, no plane) {*}{*}and no metal component.{*}{*} The carrier's ESP32-S3 PCB antenna is directly below; a ground plane 11 mm above detunes it. Wi-Fi is a supported control path per H1-Q8, so this zone is live, not void \textbar{}\\
\textbar{} {*}{*}Recovery header{*}{*} \textbar{} {*}{*}(76, 0) - (98, 20){*}{*} \textbar{} Keep clear enough that a 6-way jumper lead can be attached with this board fitted \textbar{}
\end{quote}
Note the interaction: the \textbf{DC-DC hot zone forbids aluminium electrolytics} exactly where a large bank would otherwise like to live, and the \textbf{notch removes 780 mm2} from the top edge, and the \textbf{antenna column removes a 12 x 30 mm strip} of all-layer copper from the right edge. The 2,800 uF / 100 V bank must be placed around all three.

Reverse-insertion proofing also relies on: different position counts at fixed asymmetric coordinates (a 2x7 socket cannot mate a 2x12 header); 180-degree rotation aligning neither connector; the 5th mounting hole breaking the hole pattern's rotational symmetry; and silkscreen - \textbf{a pin-1 triangle at position 1 of both blocks, plus a `\textasciicircum{}\textasciicircum{} RJ45` edge arrow} matching the carrier's.

\paragraph*{5.3 Height}
\begin{itemize}
\item Bottom side: the 8.5 mm socket bodies inside the 11.0 mm stack. \texttt{ASSUMED:} the two sockets
are the only bottom-side parts; anything else must clear the carrier's top-side components
in the same footprint region.
\item Top side: \textbf{unknown.} The 100 V bank is the tallest thing on this board and radial parts in
this class are tens of mm tall. See open question 5.
\end{itemize}
---

\subsubsection*{6. Quantity and budget}
\begin{itemize}
\item \textbf{4-6 strobe daughter boards}, of 8-12 total fixtures in the first build.
\item Also in the program: the universal carrier (8-12), the RGBW par daughter (6-8). {*}{*}The UV bar
is out of scope - D-03 closed as "no board"; allocate it no connector budget, no PWM
channels, no power budget, no enclosure provision.{*}{*}
\item System budget: \textbf{\$500-1000} for fixtures + switch + enclosures (\texttt{00} section 1). A managed
8-16 port PoE switch takes a substantial share of that, so the per-fixture allowance is
tight once the carrier, the daughter, the LED and its heatsink are counted.
\item The ICD references a \textbf{"\$30/board target"} (s1.1, in rejecting a \$15.30 mated connector
pair). Whether that figure is the carrier's or applies to this board is not stated. See
open question 7.
\end{itemize}
---

\subsubsection*{7. Assembly}
\begin{itemize}
\item Vendor: \textbf{JLCPCB PCBA}.
\item \texttt{ASSUMED:} \textbf{single-sided (top) SMD assembly}, plus THT for the two connectors.
\item The complication: the ICD requires the sockets to \textbf{face downward} - "a reverse-mounted THT
part on the daughter's bottom side, or a bottom-side SMD equivalent." Reverse-mounted THT is
not a standard JLC PCBA process, and a bottom-side SMD socket would make the assembly
double-sided. \textbf{Open question 6.}
\item Related: if the bank is built from radial aluminium electrolytics, those are THT parts too,
which lands in the same process question.
\item 4-6 boards is a hand-finishable quantity, so hand-soldering two connectors per board is a
live option rather than a fallback.
\end{itemize}
---

\subsubsection*{8. Compliance and safety flags}
Five apply. None can be guessed; the ones that need a human answer are in section 9.

\paragraph*{8.1 48 V domain - exceeds 30 V}
\begin{itemize}
\item Worst case \textbf{57 V DC} (IEEE 802.3 PSE maximum). Below the IEC 62368-1 \textbf{ES1} limit of
60 V DC, so \textbf{functional insulation only} - no basic/supplementary/reinforced safeguard,
and no safety-mandated creepage.
\item {*}{*}0.60 mm minimum outer-layer copper-to-copper spacing around every 48 V net, board-wide,
from the connector pads to the cap bank{*}{*} (IPC-2221B Table 6-1 column B2, 51-100 V band).
\textbf{This is NOT inherited automatically - this board's DRC must be set up for it.}
\item Inner layers: IPC gives 0.10 mm, below JLC's 0.127 mm minimum, so the fab minimum dominates
and the HV requirement is free there. The clearance applies \textbf{through the board too} - an
inner-layer or opposite-face signal passing under a 48 V antipad needs the same.
\item \textbf{The 0.13 mm "permanent polymer coating" column (B4) is NOT claimed.} Standard LPI
soldermask is not a qualified conformal coating, and \texttt{check\_creepage.py} implements only the
uncoated columns - a layout designed to 0.13 mm fails P8 with no waiver mechanism.
\item \textbf{Any resistor across the 48 V domain must be 0805 or larger} (0402/0603 are typically
50-75 V working) or split into two in series. {*}{*}This bites the mandatory bleed resistor and
the 48 V rail-sense divider{*}{*} - which also has to present \textless{}= 10 kohm to the carrier ADC.
\item \textbf{Capacitors on the 48 V domain must be 100 V rated.}
\end{itemize}
\paragraph*{8.2 High pulse current - \textasciitilde{}2.6 A LED drive}
Series-string drive at \textasciitilde{}2.6 A peak from a \textasciitilde{}48 V bank (STR-REQ-12). Below the 3 A continuous threshold, but pulsed and with deliberately fast edges, so it is flagged: trace width and copper weight, gate drive, current-limit accuracy and switching loop area all have to be sized against the pulse, not the average. The superseded 12 V topology would have been 8.3 A; avoiding that is the whole point of D-02.

\textbf{STR-REQ-13} - verify the candidate LED's \textbf{pulsed} forward current derating curve, not its DC maximum. The operating point is short high-current pulses at low duty cycle.

\paragraph*{8.3 Stored energy hazard - up to \textasciitilde{}3.23 J}
\begin{itemize}
\item \textasciitilde{}2,800 uF charged to 48 V holds \textbf{3.23 J}, and it holds it \textbf{with the cable unplugged}.
\item \textbf{Bleed path is mandatory} (STR-REQ-10, CAR-REQ-17), and the carrier fits no series diode
on \texttt{+48V\_SW}, so this board's bleed path is not stranded above the carrier's.
\item The board must be \textbf{safe to handle during assembly and service}. Silkscreen warning and a
defined bleed time constant are part of the deliverable.
\item A short across a charged bank is a large, fast fault current independent of any rail limit.
\end{itemize}
\paragraph*{8.4 Non-isolated topology floating at PoE potential}
\textbf{The entire fixture is non-isolated and floats at PoE potential.} Every one of these is load-bearing and lands on this board as much as on the carrier:

\begin{itemize}
\item \textbf{Non-conductive enclosure}, no chassis earth, no earthed mounting hardware bonded to
board GND.
\item {*}{*}Ethernet is the only external connection. This board may not add an external connector of
any kind.{*}{*}
\item \textbf{This board, its drive stage, its LED string and its LED wiring are all at PoE potential.}
If the LED module sits on a separate heatsink, {*}{*}that heatsink and its wiring are at PoE
potential too.{*}{*} The heatsink must not be user-accessible (H1-Q5) and must not share a mount
with anything earthed - if it is touchable, metal and earthed, the non-isolated topology is
non-conformant. (ICD notes this as carrier \texttt{decisions.md} OPEN-C, unresolved there;
H1-Q5 closes it for this program as "plastic enclosure, heatsink enclosed".)
\item \textbf{Bench hazard.} An earthed scope probe or a non-isolated USB-UART adapter ties the
floating PoE return to earth. Beyond shock and damage risk, the resulting ground currents
\textbf{break PD signature detection outright} (detection currents are only a few hundred
microamps). {*}{*}Every test point on this board carries the same silkscreen warning the
carrier's recovery header does.{*}{*}
\item \textbf{No MOV-to-earth surge network.} An unearthed PD does not need one; do not copy one out of
a reference design.
\end{itemize}
\paragraph*{8.5 Thermal - independent of firmware}
\textbf{STR-REQ-20} - over-temperature must shut the output down and assert FAULT {*}{*}without the MCU and without the network{*}{*}. \textasciitilde{}8.5 W sustained inside a sealed plastic box at 56-69 C internal air, with the LED junction also seeing peak pulses, is the case that has to fail safe on its own.

\paragraph*{8.6 Not applicable}
No mains voltage. No battery of any chemistry, no charging circuit. No motors. No RF transmitter on this board (the radio is on the carrier; this board's only RF obligation is the antenna column keepout in section 5.2).

---

\subsubsection*{9. Open questions}
Nine. Questions 1-4 change what gets designed; 5-9 are narrower but still block a decision. Defaults are offered where a sensible one exists - answering "default" to any of them is a valid answer.

---

\textbf{1. White-only or RGBW? (D-04 - this board owns it)}

Does the strobe fire only white, or can it also fire saturated colour?

\begin{itemize}
\item \textbf{White-only} covers the dominant use: rage trap and hard French rap both call for white
blasts, and SYS-REQ-07 wants clean neutral white rather than RGB-mixed white. Needs 1-2 PWM
channels and one drive stage.
\item \textbf{RGBW} adds the coloured blasts P8 (UK bass) asks for - "all fixtures fire red or green at
80-100\%". Needs 4-5 PWM channels, four drive stages, four times the LED selection and
thermal work, and it splits the same 0.99 J bank four ways.
\end{itemize}
\textbf{Recommendation: white-only.} \textbf{The one tradeoff that matters:} the 6-8 RGBW pars already produce coloured light, but they cannot strobe as hard as this board can - so choosing white-only means coloured \emph{blasts} are as fast as the pars allow, not as fast as the strobe allows. If the coloured-blast moments in P8 are meant to be percussive rather than atmospheric, that is the argument for RGBW.

\emph{Default if you have no strong feeling: white-only.}

---

{*}{*}2. How much light per flash, and what is the fastest sustained flash rate the fixture promises?{*}{*}

Currently assumed 1 J per flash at an unbounded rate. \textbf{Unbounded is not achievable} - the rail delivers 8.5 W and nothing stores more than the bank holds. The arithmetic (ICD s6.4):

\begin{itemize}
\item Bank stores \textbf{0.99 J} per full flash (48 -\textgreater{} 40 V).
\item \textbf{8.6 Hz} is the fastest rate at which every flash can be a full 0.99 J flash.
\item Above that, energy per flash must fall. At the 25 Hz ceiling (SYS-REQ-03) it is \textbf{0.34 J}.
\item A PoE+ switch later (the D-01 upgrade, resistor change only) roughly doubles both: 18.8 Hz
full-energy, 0.74 J at 25 Hz.
\end{itemize}
\textbf{And a second half to this question that the briefs have not reconciled:} STR-REQ-01 asks for a \textbf{100-200 ms} full-output flash, but 0.99 J spread over 150 ms is only about \textbf{15 W} of drive - the "100 W strobe" figure is a \textbf{10 ms} operating point, and at the closed 2.6 A drive current the bank holds full output for only \textasciitilde{}8.6 ms. So:

\begin{itemize}
\item Short accents (10-50 ms, STR-REQ-03) can be genuinely violent: \textasciitilde{}20-99 W.
\item Long flashes (100-200 ms, STR-REQ-01) land at \textasciitilde{}7-15 W.
\end{itemize}
{*}Default: accept it. Per-flash energy capped at 0.99 J; full-energy flashes up to 8.6 Hz, tapering to 0.34 J at 25 Hz, enforced by the firmware governor; "full output" for a 150 ms flash means \textasciitilde{}15 W of LED drive, and the LED is chosen so that this reads as bright in a 5 x 7 m room with 4-6 strobes firing together.{*}

{*}If you instead want a 100-200 ms flash that is blinding at 100 W, the bank has to grow by roughly 10x, which means re-opening D-02's sizing with the carrier owner - a blocking issue against LUM-CAR-A, not something this board can decide.{*}

---

\textbf{3. LED family and beam angle}

No candidate has been evaluated (STR-REQ-13 to STR-REQ-16). Two parts:

\begin{itemize}
\item \textbf{(a) Vendor/family preference?} The project plan names CREE and Lumileds as a direction
only. {*}Default: no preference - the architect picks on availability at JLCPCB, pulsed
derating data, and neutral/cool white with no colour cast at full current (a green or pink
tint at peak is a defect, STR-REQ-14).{*}
\item \textbf{(b) Beam angle?} The room is 5 x 7 m with a \textbf{2.5 m ceiling} - low. A narrow beam lights
one square metre of floor and fails the room regardless of its lumen figure (STR-REQ-16).
{*}Default: a wide beam, roughly 60-90 degrees, with an off-the-shelf lens or reflector rather
than a custom optic.{*}
\end{itemize}
---

\textbf{4. Is the LED on this board, or on a separate heatsink wired to it?}

This changes the thermal design, the board area, and whether an internal wire-to-board connection is needed. (Internal wiring is permitted - the ICD only forbids connectors that leave the enclosure.)

\begin{itemize}
\item \textbf{On-board}: emitters soldered to this PCB with a heatsink bolted through it. Fewer parts,
no wiring - but \textasciitilde{}8.5 W sustained into a 1.6 mm FR4 board sandwiched 11 mm above the carrier
in a sealed box is a poor thermal path, and it competes for area with the cap bank.
\item \textbf{Off-board}: an LED module on its own heatsink, connected by two internal wires.
\end{itemize}
\textbf{Recommendation: off-board module on its own heatsink} - it is the arrangement the ICD already assumes in section 9, and it decouples the LED's thermal problem from a crowded board. Either way the heatsink is at PoE potential and must not be user-accessible.

{*}Default: off-board, connected by a 2-pin internal wire-to-board connector or solder pads (architect's choice).{*}

---

\textbf{5. How much room is there above this board inside the enclosure?}

The 2,800 uF / 100 V bank is the tallest part on the board and parts in that class run to tens of millimetres. The enclosure is not yet designed, so this is a number that needs setting rather than discovering - and once P5 fixes the outline there is no shrink step.

\emph{Default: state a ceiling of }\emph{30 mm}{*} above the board's top face. If the intended enclosure is shallower, say so now - it forces the bank toward a larger number of shorter cans or a different capacitor technology, which changes the BOM cost and the board area.{*}

---

\textbf{6. Does the downward-facing socket force double-sided assembly at JLCPCB?}

The ICD requires the two sockets on the \textbf{bottom} side facing down. Three ways to build it:

\begin{itemize}
\item \textbf{(a)} Top-side SMD assembly at JLC, and hand-solder the two reverse-mounted THT sockets
(and any THT electrolytics) locally. Cheapest, and 4-6 boards is a hand-finishable quantity.
\item \textbf{(b)} Bottom-side SMD sockets, i.e. genuine double-sided assembly at JLC. Higher setup
cost, but nothing to hand-solder.
\item \textbf{(c)} Full JLC assembly including THT service.
\end{itemize}
\emph{Default: (a).} {*}Answer changes: the board's assembly-side constraint from P4 onward, and the capacitor technology (radial THT vs SMD).{*}

---

\textbf{7. What is the BOM cost target for one strobe daughter board, at quantity 6?}

The system budget is \$500-1000 for 8-12 fixtures plus a managed PoE switch plus enclosures. Once the carrier, the daughter, the LED and its heatsink are counted, the per-fixture allowance is tight, and the LED plus the 100 V bank are the two expensive items on this board. The ICD mentions a "\$30/board target" but does not say whether that is the carrier's figure.

\emph{Default: }\emph{\$25 per board}{*} for the PCB and its BOM at qty 6, \textbf{excluding} the LED module and heatsink, which are budgeted separately with the fixture.{*}

---

{*}{*}8. What board-type ID code is allocated to LUM-STR-A, and does this board carry an I2C EEPROM?{*}{*}

The carrier fits the top leg of the ID divider (10 k to +3V3); this board fits the bottom leg to GND, and {*}{*}the ICD says board-type codes are allocated by the carrier owner, not chosen by daughters{*}{*}. This board cannot pick its own value. The ICD also notes the Q10 default keeps an I2C EEPROM alongside the divider, without saying whether the daughter or the carrier holds it.

{*}Default: request an allocation from the carrier owner and design to a placeholder value that must be confirmed before P8; \textbf{divider only, no EEPROM on this board} (an EEPROM buys per-unit calibration storage, which matters more for the colour-matched pars than for a white strobe).{*}

---

\textbf{9. Proceed on the ICD's provisional connector coordinates, or wait for the carrier's P6?}

ICD section 7.2's connector positions - J3 at (14, 68), J4 at (56, 68), H5 at (46, 74) - are explicitly \textbf{provisional until the end of the carrier's P6}, and the ICD says daughters are blocked on them. Everything else in the ICD is frozen.

{*}Default: proceed on the provisional coordinates, and re-check them against the re-issued ICD before this board's own P6 completes. If they have moved, the fix is a \texttt{place\_edit}, not a respin. Waiting for the carrier's P6 serialises the two runs for no design benefit.{*}

---

\subsubsection*{Traceability}
Requirement IDs referenced and where they land in this document:

\begin{quote}\small\ttfamily\raggedright
\textbar{} ID \textbar{} Section \textbar{}\\
\textbar{}---\textbar{}---\textbar{}\\
\textbar{} SYS-REQ-01..08 \textbar{} 1, 3.4 (behavioural spec, inherited via STR-REQ-01..05) \textbar{}\\
\textbar{} STR-REQ-01..05 \textbar{} 1, 3.4, 4 \textbar{}\\
\textbar{} STR-REQ-06, 07 \textbar{} 2.4, 3.4 \textbar{}\\
\textbar{} STR-REQ-08, 09, 10 \textbar{} 3.2, 8.3 \textbar{}\\
\textbar{} STR-REQ-11 \textbar{} 1 (`ASSUMED:` acceptance number is \textless{} 1 ms optical rise/fall per Project Plan v1.0; the brief leaves it to be defined during design) \textbar{}\\
\textbar{} STR-REQ-12 \textbar{} 2.5, 8.2 \textbar{}\\
\textbar{} STR-REQ-13..16 \textbar{} 4, 8.2, open questions 3 and 4 \textbar{}\\
\textbar{} STR-REQ-17..22 \textbar{} 2.3, 2.4 \textbar{}\\
\textbar{} CAR-REQ-07, 08, 10..17 \textbar{} 2.3, 2.4, 3.2, 5 \textbar{}\\
\textbar{} D-01, D-02, D-03 \textbar{} 3.1, 3.2, 6 (all CLOSED) \textbar{}\\
\textbar{} D-04 \textbar{} open question 1 (OPEN, owned by this board) \textbar{}\\
\textbar{} MECH-01, MECH-02 \textbar{} 5 \textbar{}\\
\textbar{} ICD-01 \textbar{} 2, 5 (hard input, never redefined) \textbar{}\\
\textbar{} DOC-01 \textbar{} design document is a required deliverable at the end of this run \textbar{}\\
\textbar{} GIT-01 \textbar{} scoped commits only - never `gate.py --commit`, never `git add -A` \textbar{}
\end{quote}
\section{Architecture}
\subsection*{Decisions of record (state.json)}
{\small
\begin{longtable}{lp{5.2cm}p{7.2cm}}
\toprule
\textbf{Phase} & \textbf{What} & \textbf{Why} \\
\midrule
\endhead
P0 & Design baseline = white-only (D-04 recommendation), pending human verdict at H1 & Brief leans white-only; RGBW coloured blasts are already covered by 6-8 RGBW pars. Architecture is built white-only with an explicit RGBW delta documented, so a human 'RGBW' verdict costs a P2 revision not a restart. \\
P0 & BLOCKING-01: board\_init.py cannot cut the mandatory 30x26 mm RJ45 notch (ICD s7.6) & board\_swig.draw\_outline emits only a rectangle (optionally corner-filleted) on Edge.Cuts; no script in the pipeline writes Edge.Cuts geometry after P5, and the playbook states the P5 outline is final. Surfaced at H1 with options. \\
P1 & P1 roster: 4 component-scouts (led-emitter, drive-stage, energy-store, protection-sense) + 1 reference-design (pulsed-led-driver) + power-architect. No research-interface-spec agent. & The only standards-bound interface is the expansion connector, and it is frozen and fully specified in brief/06-connector-icd.md (pin map, levels, currents, creepage). An interface-spec agent would only re-transcribe it. \\
P1 & BLOCKING-02: STR-REQ-01 (100-200 ms flash at FULL output) is unachievable at any bank size that fits & 0.99 J usable / 0.15 s = 6.6 W. Reaching \textasciitilde{}92 W for 150 ms needs 13.8 J = 32,500-65,000 uF. energy-store scout measured 12,000 uF as 5,104 mm2 (64 pct of a 100x80 board) and 37 mm tall, and even that holds full output only 36.9 ms. Physics plus the closed 8.5 W budget, not a tooling limit. \\
P2 & Adopt connector ICD rev A2 (boards/lumina-carrier/architecture/connector-icd.md) as authoritative over the DRAFT-A copy in brief/06-connector-icd.md & Rev A2 adds s6.6 (binding bank-charging contract: current-limited soft-start, \textless{}=0.25 A af / \textless{}=0.5 A at, never \textgreater{}1.0 A for \textgreater{}162 ms, 3.2 J must land in the daughter limiter) and s8.4 (carrier guarantees 0 V at J3 when unprogrammed/crashed/brownout/reset), and raises board-wide 48 V outer-layer clearance from 0.60 to 0.635 mm. \\
P2 & 4-layer JLC04161H-3313; In1 AND In2 both solid GND (override planes\_gen 4L default) & Thermal forces it: check\_thermal 2L model floors at 55 C/W vs 4L 45 C/W, and the linear pass FET at 1.45 W in 56 C air fails P8 on 2L at any pour size. Supporting: 48 V routing needs 0.635 mm outer vs 0.127 mm inner, and the pulse loop return drops 1.53 mm -\textgreater{} 0.21 mm. \\
P2 & Operating point: 38.0 V string at 2.6 A, 39.7 V window floor, bank ceiling 44.5 V normal / 48.0 V armed & 38 V string reproduces the ICD s6.4 headline 0.99 J exactly while cutting board dissipation 2.33 -\textgreater{} 1.89 W (-19 pct) vs the 35.3 V alternative. The 44.5 V ceiling takes the pass FET from 1.03x its P8 thermal allowance to 1.77x margin at identical light output. \\
P2 & Charge limit 0.20 A flat current limit, not the ICD s6.6 ceiling of 0.25 A & 0.25 A x 48 V = 12.0 W instantaneous; plus the carrier's 2.4 W overhead that exceeds the 12.95 W 802.3af class envelope for far longer than the PSE 50-75 ms overload timer. 0.20 A satisfies s6.6's current-limited requirement literally and stays inside the class. \\
P2 & LED-short single-fault answer: ratiometric Vds comparator (drain vs 0.45x bank) + NTC in the pass FET's own pour, latched until ENABLE cycles & The current loop otherwise regulates 2.6 A into a shorted string, burning 125-148 W in the FET with no self-termination. A max-on-time one-shot was rejected as primary: it bounds one pulse but not the repetitive 9.9 W mean. \\
P2 & BLOCKING-03: 802.3at upgrade does not close thermally on this daughter & 4.67 W across two linear FETs in 69 C air; best case 1.91 W per D2PAK against a 1.40 W allowance = fails 1.5x. The at governor cap is \textasciitilde{}12.1 W of 18.5 W = +45 pct light, not the +120 pct D-01 implies. A real at build needs an off-board or heatsinked pass element, i.e. a respin. D-01's 'resistor change, no respin' promise does not hold for LUM-DTR-STROBE-A. \\
\bottomrule
\end{longtable}
}
\subsection*{Sheet plan}
\subsection*{LUM-DTR-STROBE-A - hierarchical sheet plan}
Four child sheets under a thin stitching root. \textasciitilde{}65 components across four cleanly separable domains, so hierarchy buys parallel P4 sheet generation and keeps the 48 V energy store visually and electrically isolated from the analogue loop and the ICD boundary.

\textbf{The sheet NAMES below are contractual.} They appear verbatim in netlist net names (\texttt{/charge/CHG\_GATE}) and \texttt{architecture/constraints.json} already references sheet-crossing names; renaming a sheet renames nets and breaks \texttt{netlist\_audit --constraints}.

\begin{quote}\small\ttfamily\raggedright
\textbar{} Sheet \textbar{} File \textbar{} Blocks \textbar{} Interface nets (sheet pins) \textbar{} refdes range \textbar{} `pwr\_base` \textbar{}\\
\textbar{}---\textbar{}---\textbar{}---\textbar{}---\textbar{}---\textbar{}---\textbar{}\\
\textbar{} root `lumina-strobe` \textbar{} `lumina-strobe.kicad\_sch` \textbar{} none (stitching only) \textbar{} - \textbar{} - \textbar{} 1 (unused) \textbar{}\\
\textbar{} {*}{*}`conn`{*}{*} \textbar{} `conn.kicad\_sch` \textbar{} J3 POWER, J4 SIGNAL, ENABLE pull-down, ID divider bottom leg, ADC RC filters, rail decoupling, mounting holes, test points \textbar{} `ENABLE`, `FAULT`, `PWM0`, `PWM1`, `PWM4`, `ADC0`, `ADC1`, `ID\_ADC`, `I2C\_SCL`, `I2C\_SDA`, `VBANK\_SENSE` \textbar{} {*}{*}1-99{*}{*} \textbar{} {*}{*}100{*}{*} \textbar{}\\
\textbar{} {*}{*}`charge`{*}{*} \textbar{} `charge.kicad\_sch` \textbar{} input TVS, hot-swap controller, charge FET, sense shunt, 2720 uF bank, passive + active bleed, bank divider \textbar{} `ENABLE`, `VBANK`, `VBANK\_SENSE`, `CHG\_EN\_n` \textbar{} {*}{*}100-199{*}{*} \textbar{} {*}{*}200{*}{*} \textbar{}\\
\textbar{} {*}{*}`drive`{*}{*} \textbar{} `drive.kicad\_sch` \textbar{} error amp, setpoint RC + SPDT, pass FET, current shunt, gate clamp, drain TVS, Vds divider, harness connector \textbar{} `VBANK`, `PWM0`, `PWM4`, `ENABLE`, `OT\_TRIP`, `UVLO\_n`, `VDS\_SENSE` \textbar{} {*}{*}200-299{*}{*} \textbar{} {*}{*}300{*}{*} \textbar{}\\
\textbar{} {*}{*}`protect`{*}{*} \textbar{} `protect.kicad\_sch` \textbar{} 2x dual comparator on `+12V` (board OT, LED OT, Vds fault latch, bank UVLO + ceiling), NTC dividers, TMP112, off-board NTC landing \textbar{} `VBANK\_SENSE`, `VDS\_SENSE`, `ENABLE`, `PWM1`, `OT\_TRIP`, `UVLO\_n`, `CHG\_EN\_n`, `FAULT`, `ADC1`, `I2C\_SCL`, `I2C\_SDA` \textbar{} {*}{*}300-399{*}{*} \textbar{} {*}{*}400{*}{*} \textbar{}
\end{quote}
\textbf{Exception to the refdes ranges:} \texttt{J3}, \texttt{J4} and \texttt{H5} keep the designators the ICD assigns them (ICD s7.2). Everything else follows the decade scheme, which is what keeps refdes unique across sheets without coordination.

---

\subsubsection*{1. Net naming - the CANONICAL final netlist names}
Mechanism (verified in this repo at \texttt{tests/s7\_regen/hierdemo/kicad/hierdemo.net}):

1. \textbf{Power SYMBOLS make a net global across the whole hierarchy} with a \textbf{bare} name and need no sheet pin. 2. A child net merged with the root through a sheet pin takes the \textbf{root-side label}: \texttt{Project.add\_sheet(child, ..., nets={[}"X"{]})} wires each sheet pin outward to a root local label \texttt{X}, so the net becomes \textbf{`/X`}. 3. A child-internal label becomes \textbf{`/\textless{}sheet\textgreater{}/NAME`}.

\textbf{These are the names `constraints.json` uses and they are contractual.} Getting one wrong makes \texttt{check\_creepage} hard-error (it raises on a declared voltage net that is not on the board) and makes \texttt{check\_current} / \texttt{check\_thermal} silently no-op.

\paragraph*{1.1 Global rails - power symbols, bare names}
\begin{quote}\small\ttfamily\raggedright
\textbar{} Net \textbar{} Present on \textbar{} Driven by \textbar{} PWR\_FLAG? \textbar{}\\
\textbar{}---\textbar{}---\textbar{}---\textbar{}---\textbar{}\\
\textbar{} `GND` \textbar{} all four \textbar{} J3 pins 2/4/6/7/8/10/13 (passive) \textbar{} {*}{*}yes{*}{*}, on `conn` \textbar{}\\
\textbar{} `+48V\_SW` \textbar{} conn, charge \textbar{} J3 pins 1/3/5 (passive) \textbar{} {*}{*}yes{*}{*}, on `conn` \textbar{}\\
\textbar{} `+12V` \textbar{} conn, drive, protect \textbar{} J3 pins 9/11 (passive) \textbar{} {*}{*}yes{*}{*}, on `conn` \textbar{}\\
\textbar{} `+3V3` \textbar{} conn, protect \textbar{} J3 pins 12/14 (passive) \textbar{} {*}{*}yes{*}{*}, on `conn` \textbar{}
\end{quote}
All four PWR\_FLAGs live on \texttt{conn}, the sheet that owns the connector, at \texttt{\#FLG0100+}. Power symbols being global, one flag each drives the net hierarchy-wide.

\textbf{`+48V\_SW` deliberately does not appear on `drive` or `protect`.} The only 48 V-domain net that crosses into \texttt{drive} is \texttt{/VBANK}, and it does so through the harness connector and the pass FET. If \texttt{+48V\_SW} ever appears on \texttt{drive} or \texttt{protect}, that is a design error, not a routing convenience - it would be a path that could energise something ahead of the charge FET.

\paragraph*{1.2 Root-crossed signals - \texttt{/NAME}}
\begin{quote}\small\ttfamily\raggedright
\textbar{} Net \textbar{} Direction \textbar{} Notes \textbar{}\\
\textbar{}---\textbar{}---\textbar{}---\textbar{}\\
\textbar{} `/ENABLE` \textbar{} conn -\textgreater{} charge, drive, protect \textbar{} Active HIGH. {*}{*}100 k pull-down on `conn`{*}{*} (ICD s8.2, mandatory). Gates the charge path's EN, the pass-FET gate clamp, and the active bleed's disarm. {*}{*}Never latched locally{*}{*} \textbar{}\\
\textbar{} `/FAULT` \textbar{} protect -\textgreater{} conn \textbar{} {*}{*}Open drain, active low.{*}{*} No pull-up on this board - the carrier's 10 k owns it. {*}{*}Never driven high{*}{*} \textbar{}\\
\textbar{} `/PWM0` \textbar{} conn -\textgreater{} drive \textbar{} {*}{*}FLASH\_GATE.{*}{*} A 5-200 ms one-shot at 1-25 Hz, LEDC timer 0. See `blocks.md` s4 \textbar{}\\
\textbar{} `/PWM1` \textbar{} conn -\textgreater{} protect \textbar{} {*}{*}BANK\_ARM.{*}{*} Static DC level, 0 \% = 44.5 V ceiling, 100 \% = 48.0 V. Timer 0, frequency-agnostic \textbar{}\\
\textbar{} `/PWM4` \textbar{} conn -\textgreater{} drive \textbar{} {*}{*}AMP\_SET.{*}{*} 13-bit at 9.766 kHz, LEDC timer 1, RC-filtered into the regulator reference \textbar{}\\
\textbar{} `/ADC0` \textbar{} conn (local) \textbar{} Bank voltage telemetry, from `/VBANK\_SENSE` through an RC at the pin. Rth 9.43 k, inside the ICD's 10 kohm \textbar{}\\
\textbar{} `/ADC1` \textbar{} protect -\textgreater{} conn \textbar{} LED-module thermistor telemetry, independent `+3V3`-referenced leg \textbar{}\\
\textbar{} `/ID\_ADC` \textbar{} conn (local) \textbar{} Board-type ID, {*}{*}bottom leg only{*}{*} - the carrier fits the 10 k top leg. {*}{*}Placeholder value; the code is allocated by the carrier owner and must be confirmed before P8{*}{*} \textbar{}\\
\textbar{} `/I2C\_SCL`, `/I2C\_SDA` \textbar{} protect \textless{}-\textgreater{} conn \textbar{} TMP112. {*}{*}No pull-ups on this board{*}{*} - the carrier's 4.7 k own the bus \textbar{}\\
\textbar{} `/VBANK` \textbar{} charge -\textgreater{} drive \textbar{} Bank positive. {*}{*}57 V, 2.6 A pulse{*}{*} \textbar{}\\
\textbar{} `/VBANK\_SENSE` \textbar{} charge -\textgreater{} conn, protect \textbar{} Divided bank voltage (2 x 82 k + 10 k). Three consumers: `ADC0`, the UVLO/ceiling comparator, the Vds comparator's reference leg \textbar{}\\
\textbar{} `/VDS\_SENSE` \textbar{} drive -\textgreater{} protect \textbar{} Divided pass-FET drain voltage (2 x 82 k + 10 k, identical ratio) \textbar{}\\
\textbar{} `/OT\_TRIP` \textbar{} protect -\textgreater{} drive, conn \textbar{} Wire-OR of three open-collector outputs (board OT, LED OT, Vds fault). Pulls the gate clamp {*}{*}and{*}{*} `FAULT` \textbar{}\\
\textbar{} `/UVLO\_n` \textbar{} protect -\textgreater{} drive \textbar{} Bank undervoltage lockout. Inhibits the drive stage only - it must {*}{*}not{*}{*} assert `FAULT` (an empty bank at power-up is not a fault) \textbar{}\\
\textbar{} `/CHG\_EN\_n` \textbar{} protect -\textgreater{} charge \textbar{} Bank-ceiling comparator output, pulls the hot-swap controller's `EN` low above the ceiling \textbar{}
\end{quote}
\paragraph*{1.3 Sheet-internal nets that \texttt{constraints.json} names}
\begin{quote}\small\ttfamily\raggedright
\textbar{} Net \textbar{} Sheet \textbar{} Why it is declared \textbar{}\\
\textbar{}---\textbar{}---\textbar{}---\textbar{}\\
\textbar{} `/charge/CHG\_GATE` \textbar{} charge \textbar{} {*}{*}64 V - the highest-voltage net on the board{*}{*} (V\_GATE-OUT is 12-16 V above VCC). Declared in `voltages` so the clearance rule covers it \textbar{}\\
\textbar{} `/drive/LED\_K` \textbar{} drive \textbar{} Harness return = pass-FET drain. {*}{*}57 V, 2.6 A{*}{*}, and the thermal-pour net for Q200 \textbar{}\\
\textbar{} `/drive/ISNS` \textbar{} drive \textbar{} Pass-FET source to shunt high side. {*}{*}2.6 A{*}{*}, but a low-voltage node - deliberately {*}{*}not{*}{*} in `voltages`, so `check\_creepage` treats it as 0 V and demands the full clearance from `/drive/LED\_K` across the FET \textbar{}
\end{quote}
\paragraph*{1.4 A merge decision worth recording}
{*}{*}\texttt{research/power.json} declares \texttt{/VBANK} and \texttt{/LED\_A} as separate power constraint entries. They are one net here, \texttt{/VBANK}.{*}{*} Nothing sits between the bank's positive terminal and the harness connector's pin 1 - no fuse, no series element - so \texttt{/LED\_A} would be an alias, and a declared net that does not exist in the netlist makes \texttt{check\_creepage} hard-error at P8. {*}{*}\texttt{/LED\_A} is deleted, not renamed.{*}{*} \texttt{/LED\_K} survives as \texttt{/drive/LED\_K} because there genuinely is a device (the string) between it and \texttt{/VBANK}.

---

\subsubsection*{2. Sheet contents and refdes allocation}
\paragraph*{2.1 \texttt{conn} - refdes 1-99, \texttt{pwr\_base} 100}
\begin{quote}\small\ttfamily\raggedright
\textbar{} Refdes \textbar{} Part \textbar{} Note \textbar{}\\
\textbar{}---\textbar{}---\textbar{}---\textbar{}\\
\textbar{} `J3` \textbar{} POWER socket, 2x7, 2.54 mm, CONNFLY DS1023-2\textbackslash{}{*}7SF11 class \textbar{} {*}{*}bottom side, faces DOWN.{*}{*} Pin map = ICD s3.1 verbatim \textbar{}\\
\textbar{} `J4` \textbar{} SIGNAL socket, 2x12, CONNFLY DS1023-2\textbackslash{}{*}12SF11 class \textbar{} {*}{*}bottom side, faces DOWN.{*}{*} Pin map = ICD s3.2 verbatim \textbar{}\\
\textbar{} `H1`-`H4` \textbar{} M3 3.2 mm mounting holes \textbar{} generated by `board\_init --mounting-holes 4`, board-only \textbar{}\\
\textbar{} `H5` \textbar{} `MountingHole\_3.2mm\_M3` symbol at (46, 74) \textbar{} {*}{*}added at P4 as a symbol so it carries a refdes{*}{*} - `--mounting-holes` makes corner holes only \textbar{}\\
\textbar{} `R1` \textbar{} 100 k 0603 \textbar{} {*}{*}ENABLE pull-down, ICD-mandated{*}{*} \textbar{}\\
\textbar{} `R2` \textbar{} ID\_ADC bottom leg, 0603 1 \% \textbar{} {*}{*}placeholder value - confirm with the carrier owner before P8{*}{*} \textbar{}\\
\textbar{} `R3`, `C3` \textbar{} ADC0 RC (1 k + 10 nF) \textbar{} 9.43 k source impedance is at the ICD ceiling; SAR sampling charge wants a local reservoir \textbar{}\\
\textbar{} `R4`, `C4` \textbar{} ADC1 RC \textbar{} same \textbar{}\\
\textbar{} `C1` \textbar{} `+12V` local bulk, {*}{*}\textless{}= 4.7 uF{*}{*} \textbar{} deliberately small so the error amp dies promptly on unplug \textbar{}\\
\textbar{} `C2`, `C5` \textbar{} `+3V3` and `+12V` 100 nF \textbar{} \textbar{}\\
\textbar{} `TP1`-`TP6` \textbar{} test points: `+48V\_SW`, `/VBANK`, `GND`, `/drive/ISNS`, `/OT\_TRIP`, `/ENABLE` \textbar{} {*}{*}every one carries the floating-PoE silkscreen warning{*}{*} \textbar{}
\end{quote}
\paragraph*{2.2 \texttt{charge} - refdes 100-199, \texttt{pwr\_base} 200}
\begin{quote}\small\ttfamily\raggedright
\textbar{} Refdes \textbar{} Part \textbar{} Note \textbar{}\\
\textbar{}---\textbar{}---\textbar{}---\textbar{}\\
\textbar{} `D100` \textbar{} TVS, SMBJ58A class, SMB \textbar{} 58 V standoff / 93.6 V clamp, under the 100 V cap rating \textbar{}\\
\textbar{} `U100` \textbar{} Hot-swap / power-limiting controller, TPS2490 class, MSOP-10 \textbar{} {*}{*}ILIM 0.20 A, PLIM 12 W, fault timer \textgreater{} 653 ms.{*}{*} EN is GND-referenced \textbar{}\\
\textbar{} `Q100` \textbar{} Charge FET, D2PAK N-channel 100 V, IRF540N class \textbar{} Tab = drain = `+48V\_SW`. {*}{*}\textgreater{}= 645 mm2 pour{*}{*} \textbar{}\\
\textbar{} `R100` \textbar{} Sense resistor, 0805 or 1206, value from the controller's threshold at P3 \textbar{} \textasciitilde{}250 mohm if the threshold is 50 mV \textbar{}\\
\textbar{} `R101`, `R102`, `R103`, `C100`, `C101`, `C102`, `R104` \textbar{} PROG/PLIM divider, TIMER, VCC bypass, dv/dt gate cap, EN series \textbar{} \textbar{}\\
\textbar{} `C110`-`C113` \textbar{} Bank HF: 4 x 10 uF / 100 V X7S 1210 \textbar{} {*}{*}2.7 uF each effective at 48 V{*}{*} \textbar{}\\
\textbar{} `C120`-`C125` \textbar{} Bank bulk: {*}{*}six D18 / 7.5 mm radial footprints, four populated{*}{*} at 680 uF / 100 V \textbar{} Second-sources the bank at 470 uF across four vendors; 2720 -\textgreater{} 4080 uF knob with no respin \textbar{}\\
\textbar{} `R110` \textbar{} Passive bleed, 100 k 0805, 150 V working \textbar{} 23 mW, 272 s. {*}{*}Un-defeatable{*}{*} \textbar{}\\
\textbar{} `R111`, `R112` \textbar{} Active bleed, 2 x 470 ohm 2512 2 W \textbar{} 1.23 W peak each, inside the continuous rating - {*}{*}no joule rating needed, and none is published{*}{*} \textbar{}\\
\textbar{} `Q110` \textbar{} Bleed switch, 150 V N-ch SOT-223, CJT04N15 class \textbar{} \textbar{}\\
\textbar{} `Q111`, `R113`, `D110` \textbar{} Disarm 2N7002, 1 M bias from `/VBANK`, 10 V zener clamp \textbar{} {*}{*}Bias UP from the bank, pull DOWN to disarm{*}{*} - so "every rail dead" is the ON state \textbar{}\\
\textbar{} `R120`, `R121`, `R122` \textbar{} Bank divider, 2 x 82 k + 10 k, 0805 \textbar{} Rth 9.43 k. Satisfies the 0805 rule {*}and{*} the series-split rule at once \textbar{}
\end{quote}
\paragraph*{2.3 \texttt{drive} - refdes 200-299, \texttt{pwr\_base} 300}
\begin{quote}\small\ttfamily\raggedright
\textbar{} Refdes \textbar{} Part \textbar{} Note \textbar{}\\
\textbar{}---\textbar{}---\textbar{}---\textbar{}\\
\textbar{} `Q200` \textbar{} Pass FET, D2PAK planar HEXFET-5, 200 V / 18 A, IRF640N class \textbar{} Tab = drain = `/drive/LED\_K`. {*}{*}\textgreater{}= 900 mm2 pour, target 1000 mm2, \textgreater{}= 12 thermal vias{*}{*} \textbar{}\\
\textbar{} `R200` \textbar{} Shunt, 200 mohm 3 W 2512 1 \% \textbar{} 520 mV FS. {*}{*}Kelvin layout{*}{*} - sense traces to the pad ends \textbar{}\\
\textbar{} `U200` \textbar{} Dual op-amp, LM2904 class, SOIC-8, {*}{*}on `+12V`{*}{*} \textbar{} A = error amp / gate driver; B = setpoint buffer, free \textbar{}\\
\textbar{} `U210` \textbar{} SPDT analogue switch, SGM3157 class, SC-70-6 \textbar{} Steers the {*}{*}reference{*}{*}, not the gate \textbar{}\\
\textbar{} `R201`, `R202`, `C200` \textbar{} Setpoint RC + divider: 10 k + 100 nF, 5k36/1k00 (6.35:1) \textbar{} tau 1 ms, 4.6 ms to 1 \%, 2.6 \% ripple at 9.766 kHz \textbar{}\\
\textbar{} `R203` \textbar{} Gate pull-down 4k7 \textbar{} {*}{*}The interlock of record{*}{*} - off with {*}everything{*} dead \textbar{}\\
\textbar{} `R204` \textbar{} Gate series \textbar{} \textbar{}\\
\textbar{} `Q210`, `Q211` \textbar{} 2 x 2N7002 (JLC {*}{*}Basic{*}{*}) \textbar{} One inverts ENABLE, one clamps gate-to-source. Also the landing point for `/OT\_TRIP` and `/UVLO\_n` \textbar{}\\
\textbar{} `D200` \textbar{} Drain-source TVS, SMBJ58A class \textbar{} Harness `L di/dt` clamp. {*}{*}NOT a freewheel diode across the string{*}{*} \textbar{}\\
\textbar{} `R210`, `R211`, `R212` \textbar{} Vds divider, 2 x 82 k + 10 k, 0805 \textbar{} Identical ratio to the bank divider, so the fault trip is ratiometric \textbar{}\\
\textbar{} `J200` \textbar{} Harness, JST VH 2-pin THT, 10 A / 250 V \textbar{} `/VBANK` and `/drive/LED\_K` to the off-board module. {*}{*}Top edge{*}{*} \textbar{}\\
\textbar{} `TP200`-`TP202` \textbar{} shunt Kelvin probe point, gate, `/drive/LED\_K` \textbar{} short-return probe point at the shunt - a standard ground lead cannot measure this edge \textbar{}
\end{quote}
\paragraph*{2.4 \texttt{protect} - refdes 300-399, \texttt{pwr\_base} 400}
\begin{quote}\small\ttfamily\raggedright
\textbar{} Refdes \textbar{} Part \textbar{} Note \textbar{}\\
\textbar{}---\textbar{}---\textbar{}---\textbar{}\\
\textbar{} `U300` \textbar{} Dual comparator, LM2903 class, SOIC-8, {*}{*}on `+12V`{*}{*} \textbar{} A = board OT, B = LED-module OT \textbar{}\\
\textbar{} `U301` \textbar{} Dual comparator, LM2903 class, SOIC-8, {*}{*}on `+12V`{*}{*} \textbar{} A = Vds fault (latched), B = bank UVLO + ceiling \textbar{}\\
\textbar{} `RT300` \textbar{} NTC 10 k 0603, {*}{*}in Q200's drain pour{*}{*} \textbar{} \textbar{}\\
\textbar{} `RT301` \textbar{} NTC 10 k 0603, {*}{*}in Q100's drain pour{*}{*} \textbar{} {*}{*}In parallel with RT300{*}{*} as one divider's bottom leg - the hotter device dominates, so one comparator section ORs both \textbar{}\\
\textbar{} `R300`-`R309` \textbar{} OT dividers and ratiometric references from `+12V` \textbar{} supply-independent trip \textbar{}\\
\textbar{} `R310`-`R315`, `C300` \textbar{} Vds threshold (\textasciitilde{}0.45), arming RC (10 k + 10 nF = 100 us blank) \textbar{} \textbar{}\\
\textbar{} `Q300` \textbar{} 2N7002, fault latch \textbar{} {*}{*}Cleared only by `ENABLE` going low{*}{*} - a fault latch, not an ENABLE latch \textbar{}\\
\textbar{} `R320`-`R323` \textbar{} Bank UVLO and ceiling thresholds, hysteresis \textbar{} UVLO trip = `V\_string + 1.7 V`, {*}{*}set by one 1 \% resistor at P3 from the measured string{*}{*} \textbar{}\\
\textbar{} `U310` \textbar{} TMP112 class, SOT-563, {*}{*}on `+3V3`{*}{*} \textbar{} I2C telemetry. {*}{*}Fit no pull-ups{*}{*} \textbar{}\\
\textbar{} `J300` \textbar{} 4-way internal landing for the two off-board thermistors \textbar{} LED-module NTC (trip, `+12V`) and LED-module NTC (telemetry, `+3V3`). {*}{*}JLC stocks no leaded/probe NTC at all{*}{*} - these are hand-terminated, not a PCBA line \textbar{}\\
\textbar{} `R330`, `R331` \textbar{} LED-module NTC divider, {*}{*}NTC as the TOP leg{*}{*} \textbar{} an open harness wire pulls the node low and {*}{*}trips{*}{*} - fail-safe on a broken wire \textbar{}
\end{quote}
---

\subsubsection*{3. What P4 must review before generating}
Three items that are invisible on the bench and expensive at P8:

1. \textbf{The 48 V domain boundary.} No component may bridge \texttt{+12V} or \texttt{+3V3} to \texttt{+48V\_SW}, \texttt{/VBANK} or \texttt{/drive/LED\_K} (ICD s8.3 point 2). The only nets crossing are \texttt{ENABLE}/\texttt{PWM} into GND-referenced inputs and the two divider taps, which run \textbf{out} of the 48 V domain into high-impedance inputs. Violating this re-creates the 802.3 port-capacitance problem behind the compliance switch's back and shows up only as a PD compliance failure. 2. \textbf{`FAULT` has no pull-up on this board and is never driven by a push-pull output.} Every device touching it is open collector or open drain. 3. \textbf{No I2C pull-ups.} The carrier's 4.7 k own the bus.

Full architecture narrative (not inlined; see the artifact index): \texttt{architecture/blocks.md}, \texttt{architecture/power\_tree.md}, \texttt{architecture/stackup.md}.
\section{Schematic}
\emph{Pending --- produced at P4.}
\section{Layout}
\emph{Pending --- produced at P6.}
\section{Verification}
\emph{Pending --- produced at P8.}
\section{DFM and Fabrication}
\emph{Pending --- produced at P9.}
\section{Run Record (Appendix)}
\subsection*{Phase digests}
\paragraph*{P0}
\begin{quote}\small\ttfamily\raggedright
\# P0 Intake - digest\\
~\\
- `requirements-analyst` read all five brief files and wrote `requirements.md` (627 lines).\\
  Both connector pin maps transcribed verbatim from the ICD; nothing invented.\\
- Closed decisions recorded as settled, not re-raised: D-01 (at power stage / af\\
  classification, budgets against af), D-02 (\textasciitilde{}2800 uF bank on the 48 V raw rail, 100 V\\
  parts, \textasciitilde{}2.6 A series string, mandatory bleed), D-03 (UV bar out of scope), the ICD rail\\
  contract, the 802.3 port-capacitance clause, the full mechanical footprint including the\\
  30 x 26 mm RJ45 notch and three keepouts, and the corrected 8.5 W (af) / 18.5 W (at)\\
  budget.\\
- Safety flags raised: 48 V domain (57 V worst case), 2.6 A pulsed drive, 3.23 J stored\\
  energy with the cable unplugged, non-isolated topology floating at PoE potential\\
  (LED wiring and heatsink included), sealed-box ambient 56-69 C.\\
- Substantive finding: STR-REQ-01 (100-200 ms flash at FULL output) and the closed bank\\
  sizing are not consistent. 0.99 J over 150 ms is \textasciitilde{}15 W of drive, not the briefs' 100 W.\\
  Folded into open question 2 rather than silently resolved.\\
- Source conflict recorded, not averaged: `05` gives 8.6-9.3 W (af) to the daughter, the\\
  ICD s6.2 gives 8.5 W total. Per the ICD's own precedence clause the ICD governs.\\
- 9 open questions produced. D-04 (white-only vs RGBW) is this board's to close.\\
- Tooling gap found while checking mechanics: `board\_init.py` cannot cut the mandatory\\
  notch (BLOCKING-01, logged as a P0 decision).
\end{quote}
\paragraph*{P1}
\begin{quote}\small\ttfamily\raggedright
\# P1 Research - digest\\
~\\
Roster: 4 `research-component-scout` (led-emitter, drive-stage, energy-store,\\
protection-sense) + 1 `research-reference-design` (pulsed-led-driver) + 1\\
`research-power-architect`. No `research-interface-spec`: the only standards-bound\\
interface is the expansion connector and it is frozen and fully specified.\\
~\\
- {*}{*}led-emitter{*}{*}: no JLC-stocked white LED is DC-rated for 2.6 A, and every published\\
  pulse rating in this market is at \textasciitilde{}100 us / 10 \% duty - 50x to 2000x shorter than the\\
  5-200 ms flashes here. Cree's XP-G2 datasheet contains the word "pulse" zero times.\\
  Design rule adopted: budget ZERO pulsed headroom, keep peak inside each die's DC max.\\
  Verdict: array goes OFF-BOARD on an aluminium MCPCB + heatsink (on-board FR4 would be\\
  120-200 C of rise). The right optic is no optic - bare 120 deg at 2.3 m already covers\\
  the 5x7 m room; any TIR narrows it to the STR-REQ-16 failure mode.\\
- {*}{*}drive-stage{*}{*}: JLC has no HV linear CC driver near 2.6 A and no wide-SOA linear-mode\\
  MOSFET at all, so the discrete op-amp + shunt + FET loop is forced, not preferred.\\
  200 mohm 2512 shunt beats 50 mohm (13 \% vs 54 \% offset error at STR-REQ-04's 10 \% dim\\
  point). Found the LED-short fault: the loop keeps regulating 2.6 A into a shorted\\
  string, 125-148 W in the FET, no self-termination.\\
- {*}{*}energy-store{*}{*}: 4 x 680 uF/100 V = 2720 uF, 27 mm tall, \$6.18/board. Ripple and ESR\\
  are NOT binding (8.7x ripple margin, 1.43 \% IR sag); height, endurance and Extended\\
  fees are. Growing the bank 4x does not fit (64 \% of the board, 37 mm tall) and even\\
  then holds full output only 36.9 ms.\\
- {*}{*}refdesign-pulsed-led-driver{*}{*}: TI TIDA-01081 is direct prior art. The decay tail is\\
  the driver's inductor, output capacitance and soft-start - never the LED (phosphor\\
  persistence \textasciitilde{}60 ns, four orders below the 1 ms spec). Shunt-FET dimming REJECTED: it\\
  keeps burning current while dark. Linear-mode SOA rules: pick older generations and\\
  higher voltage classes, never select on Rds(on).\\
- {*}{*}power{*}{*}: zero local regulators, zero magnetics, one current limiter. Board\\
  dissipation is an invariant, `P\_rail x (48 - V\_string)/48 + housekeeping` - string\\
  voltage is the only first-order lever. Exactly ONE full-energy flash from a full bank.\\
  And the finding that matters beyond this board: {*}{*}the 802.3at upgrade does not close\\
  thermally.{*}{*}
\end{quote}
\paragraph*{P2}
\begin{quote}\small\ttfamily\raggedright
\# P2 Architecture - digest\\
~\\
Adopted connector ICD {*}{*}rev A2{*}{*} (`boards/lumina-carrier/architecture/connector-icd.md`)\\
over the DRAFT-A copy in `brief/`: it adds s6.6 (binding current-limited bank-charging\\
contract) and s8.4 (carrier guarantees 0 V at J3 when unprogrammed/crashed/reset), and\\
raises board-wide 48 V outer clearance from 0.60 to 0.635 mm.\\
~\\
- {*}{*}4 layers, `JLC04161H-3313`.{*}{*} Forced by thermals, not routing: `check\_thermal`'s 2L\\
  model floors at 55 C/W vs 4L's 45 C/W, and the linear pass FET fails P8 on 2 layers at\\
  ANY pour size. In1 AND In2 both solid GND - a deliberate override of `planes\_gen`'s 4L\\
  default. Antenna column is voided on all four layers.\\
- {*}{*}Four sheets{*}{*}: conn (ICD boundary), charge (0.20 A hard-limited hot-swap + 2720 uF /\\
  100 V bank + self-powered bleed), drive (op-amp + 200 mohm shunt + D2PAK planar HEXFET,\\
  no inductor and no output cap), protect (2x dual comparator on +12V: board OT, LED OT,\\
  Vds fault latch, bank UVLO + ceiling).\\
- {*}{*}Operating point: 38.0 V string at 2.6 A, 39.7 V window floor, bank ceiling 44.5 V\\
  normal / 48.0 V armed.{*}{*} 38 V reproduces the ICD's own headline 0.99 J exactly while\\
  cutting board dissipation 2.33 -\textgreater{} 1.89 W (-19 \%). The 44.5 V ceiling is what takes the\\
  pass FET from 1.03x its P8 allowance to 1.77x margin, at identical light output.\\
- {*}{*}Charge limit 0.20 A, not the ICD's 0.25 A ceiling{*}{*}: 0.25 A is 12.0 W instantaneous\\
  which, plus the carrier's 2.4 W, exceeds the 12.95 W af class envelope for longer than\\
  the PSE's 50-75 ms overload timer. Also found: the charge path is a permanently-cycling\\
  limiter in regulation \textasciitilde{}88 \% of the time, not a start-up element.\\
- {*}{*}LED-short fault decided{*}{*} (not deferred): ratiometric Vds comparator, drain vs\\
  0.45 x bank off identical dividers so no voltage reference is needed, \textasciitilde{}20 us,\\
  arming-blanked, latched until ENABLE cycles - plus an NTC in the pass FET's own pour.\\
  A max-on-time one-shot was rejected as primary: it bounds one pulse, not the repetitive\\
  9.9 W mean.\\
- {*}{*}Emitter re-decided{*}{*} off the false JLC constraint (the array is off-board, so it is\\
  not a JLC PCBA line item): single 3S string of \textgreater{}=2.6 A-DC-rated multi-die emitters\\
  (XHP70.3 12 V class). No parallel paths means no ballast and no Vf binning. P3's\\
  part-sourcer will not find it on JLC; that is correct and expected.\\
- {*}{*}Cost \textasciitilde{}\$24.50/board at qty 6{*}{*} (\textasciitilde{}\$14.00 BOM + \textasciitilde{}\$3 PCB + \textasciitilde{}\$7.50 Extended setup),\\
  against the \$25 default target. Light engine separate at \textasciitilde{}\$46-65/fixture.\\
- Two things carried to H1: {*}{*}802.3at does not close thermally on this daughter{*}{*}, and\\
  {*}{*}the RJ45 notch cannot be cut by this pipeline{*}{*}.
\end{quote}
\subsection*{Gate attempt history}
\emph{(no entries)}
\subsection*{Issues}
\emph{(no entries)}
\subsection*{Human checkpoints}
{\small
\begin{longtable}{lllp{2.6cm}p{6.4cm}}
\toprule
\textbf{Id} & \textbf{Phase} & \textbf{Status} & \textbf{When} & \textbf{Note} \\
\midrule
\endhead
1 & P2 & not reached &  &  \\
2 & P4 & not reached &  &  \\
3 & P6 & not reached &  &  \\
4 & P8 & not reached &  &  \\
5 & P10 & not reached &  &  \\
\bottomrule
\end{longtable}
}
\section{Artifact Index}
{\small
\begin{longtable}{p{9cm}rp{4cm}}
\toprule
\textbf{File} & \textbf{Size} & \textbf{Note} \\
\midrule
\endhead
\texttt{state.json} & 8,989 B &  \\
\texttt{requirements.md} & 34,104 B &  \\
\texttt{architecture/blocks.md} & 32,529 B &  \\
\texttt{architecture/constraints.json} & 6,752 B &  \\
\texttt{architecture/power\_tree.md} & 36,997 B &  \\
\texttt{architecture/sheets.md} & 13,783 B &  \\
\texttt{architecture/stackup.md} & 19,548 B &  \\
\bottomrule
\end{longtable}
}
\end{document}
